\chapter{Afweging van alternatieve oplossingen}
\label{chap:alternatieven}

\section{Desktopapplicatie tegenover webapplicatie}
Een centrale webapplicatie was een mogelijk alternatief voor Electron. Een webserver zou updates voor alle gebruikers onmiddellijk beschikbaar maken en meerdere clients kunnen bedienen. Daartegenover staan extra operationele afhankelijkheden: hosting, authenticatie, netwerkconfiguratie en bereikbaarheid van de server. De timingbron is zelf al afhankelijk van internet. Door verwerking en opslag lokaal te houden wordt ten minste de reeds ontvangen historiek niet van een tweede externe dienst afhankelijk.

Een gewone lokale webserver met browserinterface had eveneens gekund. Electron voegt pakketgrootte en geheugengebruik toe, maar vereenvoudigt native bestandskeuze, vensterbeheer, installatie en distributie. Vooral het openen van afzonderlijke dashboards voor extra wagens en grafieken past goed bij een desktopapplicatie. Voor dit prototype woog die ontwikkelsnelheid zwaarder dan de grotere runtime.

Een volledig native Windowsapplicatie, bijvoorbeeld in C\#, zou een kleinere en nauwer geïntegreerde productieapp kunnen opleveren. De ontwikkeling gebeurde echter op macOS en het resultaat moest ook daar testbaar blijven. Electron maakte één gedeelde implementatie mogelijk. Wanneer de applicatie later zeer breed wordt uitgerold en resourcegebruik problematisch blijkt, kan een native frontend opnieuw worden onderzocht terwijl de pure domeinregels als specificatie behouden blijven.

\section{DOM-extractie tegenover een officiële API}
Een officiële JSON- of WebSocket-API zou betrouwbaarder zijn dan DOM-extractie. Velden zouden een gedocumenteerd type en vaste semantiek kunnen hebben. Voor de gebruikte publieke timingpagina's was zo'n uniforme beschikbare API niet het uitgangspunt. De applicatie leest daarom de gerenderde tabel die ook de gebruiker ziet.

Deze oplossing is pragmatisch maar kwetsbaar voor HTML-wijzigingen. De impact wordt beperkt door headercanonicalisatie, providerdetectie en parserdiagnostiek. Wanneer in de toekomst toegang tot een officiële feed beschikbaar wordt, kan een nieuwe collectoradapter worden toegevoegd. De downstreammodules hoeven dan niet te veranderen zolang dezelfde genormaliseerde rijen worden geproduceerd.

\section{Bestanden tegenover SQLite}
SQLite was aanvankelijk de voor de hand liggende keuze. Het biedt transacties, indexes en krachtige queries zonder afzonderlijke server. Voor de uiteindelijke schaal bleken bestanden echter voordelen te hebben. Een engineer kan een CSV direct openen, JSONL kan append-only worden geschreven en een beschadigde laatste regel tast niet noodzakelijk de volledige historiek aan. Ook export en debugging zijn eenvoudiger.

Het nadeel is dat relaties en constraints in applicatiecode worden bewaakt. Queries over meerdere grote evenementen zijn minder efficiënt. Wanneer later een centraal archief met vele races en interactieve historische analyses nodig is, wordt SQLite of PostgreSQL relevanter. De huidige storage schema-module biedt daarvoor al een uniforme representatie; migratie hoeft niet vanuit willekeurige rendererdata te gebeuren.

\section{Volledig incrementele tegenover herberekende statistiek}
Een gemiddelde kan efficiënt worden bijgewerkt met een teller en som. Voor een lange race lijkt dit aantrekkelijk. In dit project zijn ronden echter niet altijd definitief geldig op het eerste moment dat ze verschijnen. Pitstatus, sectorvlaggen of correcties kunnen de selectie beïnvloeden. Een gecachet gemiddelde zou dan teruggedraaid moeten worden met exact dezelfde eerdere context.

Daarom wordt de lap history als bron gebruikt en worden statistieken opnieuw uit de gefilterde relevante ronden berekend. Met enkele duizenden records en pure JavaScriptfuncties blijft dit snel genoeg. De laatst berekende samenvatting wordt wel opgeslagen voor inspectie. Pas wanneer metingen aantonen dat herberekenen een werkelijk prestatieprobleem vormt, is een incrementele cache met invalidatieregels verantwoord.

\section{Eenvoudige voorspelling tegenover machine learning}
Een regressiemodel, random forest of neuraal netwerk zou meerdere features kunnen combineren: piloot, bandenslijtage, brandstofniveau, verkeer, temperatuur en weersomstandigheden. De timingpagina levert veel van die inputs niet. De beschikbare trainingsdata is bovendien beperkt en bevat veranderende circuits en omstandigheden. Een complex model zou daardoor snel overfitten en moeilijk uitlegbaar zijn.

Het sectorgemiddeldenmodel heeft duidelijke beperkingen, maar iedere component kan op het dashboard worden verklaard. Het model reageert op actuele sectoren en gebruikt alleen recente data van dezelfde piloot. Een toekomstige uitbreiding kan deze baseline vergelijken met exponentiële weging of regressie. Een complexer model is pas zinvol wanneer het op historische sessies aantoonbaar een lagere fout heeft en zijn onzekerheid operationeel kan communiceren.

\section{Automatische tegenover handmatige modusdetectie}
De sessienaam bevat vaak woorden als Race, Qualifying of Practice. Automatische detectie lijkt dus eenvoudig. Toch verschillen talen, afkortingen en providerlabels, en een organisator kan een afwijkende naam gebruiken. Een fout gedetecteerde modus kan pitinformatie verbergen of verkeerde vergelijkingen tonen. Daarom kiest de gebruiker de modus expliciet. De extra handeling is klein en het gedrag blijft voorspelbaar.

\section{Iteratieve tegenover vooraf vastgelegde ontwikkeling}
Een watervalaanpak met een volledige specificatie vóór implementatie zou minder herwerking beloven. In deze opdracht bestond veel kennis echter impliciet bij de opdrachtgever. Pas een zichtbaar dashboard maakte duidelijk welke informatie ontbrak. Een iteratieve aanpak met korte feedbacklussen was daardoor geschikter.

Dit betekende niet dat elke nieuwe wens onmiddellijk zonder afweging werd toegevoegd. De architectuur werd gaandeweg strenger: gedeelde logica moest naar pure modules, nieuwe edgecases kregen regressietests en providerafhankelijke aannames moesten expliciet worden benoemd. Zo bleef het prototype ondanks veranderende requirements onderhoudbaar.

